\uuid{PKRL}
\titre{Dérivée}
\auteur{bodin ; hanani ; mzari}
\organisation{exo7}
\licence{CC BY-NC-SA 4.0}
\datecreate{2020-01-20}
\niveau{L1}
\module{Algèbre}
\chapitre{Polynôme, fraction rationnelle}
\sousChapitre{Définition, degré, produit}
\difficulte{2}
\theme{Polynômes, dérivée}
\exo7qcmid{lille-1.polynomes.005}

\contenu{
\begin{question}
\qclassification{105.05}
\qtags{motcle=dérivée}
Soit $P(X) = \sum_{k=0}^n a_k X^k$. On associe le polynôme dérivé :
$P'(X) = \sum_{k=1}^n ka_k X^{k-1}$.
\begin{answers}
    \bad{Le polynôme dérivé de $P(X) = X^5-2X^2+1$ est $P'(X)=5X^4-2X$.}

    \bad{Le seul polynôme qui vérifie $P'(X)=0$ est $P(X)=1$.}

    \good{Si $P'(X)$ est de degré $7$, alors $P(X)$ est de degré $8$.}

    \bad{Si le coefficient constant de $P$ est nul, alors c'est aussi le cas pour $P'$.}
\end{answers}
\begin{explanations}
Le polynôme dérivé s'obtient comme si on dérivait la fonction $X \mapsto P(X)$.
\end{explanations}
\end{question}
}
